\documentclass[11pt]{article}

% ── Packages ──────────────────────────────────────────────────────────────────
\usepackage{amsmath, amsthm, amssymb}
\usepackage[colorlinks=true, linkcolor=blue, citecolor=blue, urlcolor=blue]{hyperref}
\usepackage{booktabs}
\usepackage{graphicx}
\usepackage{xcolor}
\usepackage{microtype}
\usepackage[margin=1in]{geometry}

% ── Theorem environments ───────────────────────────────────────────────────────
\newtheorem{theorem}{Theorem}[section]
\newtheorem{conjecture}[theorem]{Conjecture}
\newtheorem{proposition}[theorem]{Proposition}
\newtheorem{definition}[theorem]{Definition}
\newtheorem*{openprob}{Open Problem}

\theoremstyle{remark}
\newtheorem*{remark}{Remark}

% ── Macros ─────────────────────────────────────────────────────────────────────
\newcommand{\qubo}{QUBO\xspace}
\newcommand{\qaoa}{QAOA\xspace}
\newcommand{\calA}{\mathcal{A}}
\newcommand{\calB}{\mathcal{B}}
\newcommand{\calC}{\mathcal{C}}
\newcommand{\Var}{\operatorname{Var}}
\newcommand{\E}{\mathbb{E}}

% ── Title ──────────────────────────────────────────────────────────────────────
\title{\textbf{Encoding Complexity, Barren Plateaus, and the\\
Circularity Conjecture: Structural Limits on\\
Hybrid Quantum-Classical Optimisation}}

\author{[Author Name]\\
\small Independent Researcher\\
\small \href{https://github.com/jasstt/quantum-knapsack-vrp-theory}%
{\texttt{github.com/jasstt/quantum-knapsack-vrp-theory}}}

\date{\today}

% ══════════════════════════════════════════════════════════════════════════════
\begin{document}
\maketitle

% ── Abstract ──────────────────────────────────────────────────────────────────
\begin{abstract}
The Quantum Approximate Optimization Algorithm (QAOA) has generated
substantial interest as a near-term approach to combinatorial optimisation,
yet empirical results have consistently fallen short of theoretical promise.
We investigate the structural conditions under which quantum advantage is
theoretically accessible, combining an analytical QUBO encoding complexity
classification across seven problem classes with an empirical measurement of
barren plateau gradient variance scaling on MaxCut QAOA circuits.
We classify problems by whether their QUBO formulation grows as $O(n)$,
$O(nk)$, or $O(n^2m)$ in the number of binary variables, and show that this
growth rate directly controls the effective barren plateau suppression---from
polynomial ($O(n)$ class) to doubly-exponential ($O(n^2m)$ class, including
VRP and TSP).
For the 20-node CVRP benchmark, our encoding analysis predicts
$\approx 2{,}092$ logical and $\approx 939{,}000$ physical qubits, consistent
with empirical resource estimates of 2,648 logical and 1.2 million physical
qubits.
We introduce the \emph{Circularity Conjecture}: for natural combinatorial
problem classes, the problems where LP boundary zones are small enough for
quantum sub-problems to be tractable are precisely those where classical
LP-plus-rounding already provides strong approximation guarantees.
We further identify as an open problem the question of tight lower bounds on
subtour encoding size for TSP and VRP.
This structural coincidence is not a hardware limitation and will not be
resolved by improved quantum devices alone---it suggests that the search for
practical quantum advantage in combinatorial optimisation must target problem
classes beyond those studied in this work.
\end{abstract}

\tableofcontents
\newpage

% ══════════════════════════════════════════════════════════════════════════════
\section{Introduction}
\label{sec:intro}

The Quantum Approximate Optimization Algorithm (QAOA) generated substantial interest upon its introduction in 2014~\cite{farhi2014qaoa}. It was presented as the most promising candidate for near-term quantum advantage in combinatorial optimisation. However, over the past decade, empirical results have consistently fallen short of theoretical promises---even for MaxCut, QAOA has not surpassed the classical 0.878 approximation guarantee of the Goemans--Williamson algorithm. This study investigates the underlying reasons for this shortfall, asking the ``why'' question from a structural perspective rather than a purely algorithmic or hardware-focused one.

We present two parallel contributions. Analytically, we classify seven combinatorial problem classes by deriving the growth rate of their QUBO encoding size relative to the instance size $n$, ranging from $O(n)$ to $O(n^2m)$. Empirically, we measure the barren plateau gradient variance on MaxCut QAOA circuits, verifying the $O(2^{-n})$ scaling established by McClean et al.~\cite{mcclean2018barren}. By linking these two perspectives, we investigate how the encoding size directly controls the severity of the barren plateau.

We introduce the Circularity Conjecture: the combinatorial problems where quantum sub-problems are small enough to remain tractable (characterised by $O(n)$ encoding, local cost functions, and high valid solution fractions) are exactly the problems where classical linear programming and rounding techniques already work well. This coincidence constitutes a structural barrier, not a hardware limitation, suggesting it will not be resolved simply by building better quantum devices. These findings are supported by empirical data from a 20-node CVRP benchmark, which requires approximately 1.2 million physical qubits, consistent with its theoretical $O(n^2m)$ classification.

The structure of this paper is as follows. Section~\ref{sec:encoding} derives the QUBO encoding complexity classes. Section~\ref{sec:barren} details the connection between encoding complexity and barren plateaus. Section~\ref{sec:circularity} formalises the Circularity Conjecture and explores potential proof directions. Section~\ref{sec:exceptions} examines candidate exceptions to this circularity, such as MaxCut and the $k$-Local Hamiltonian problem. Finally, Section~\ref{sec:vrp} interprets the VRP benchmark results in light of our theoretical framework.

% ══════════════════════════════════════════════════════════════════════════════
\section{Problem Encoding Complexity Classes}
\label{sec:encoding}

\subsection{Derivation Framework}

A QUBO formulation for a combinatorial optimisation problem requires:
\begin{equation}
  H(x) = \underbrace{-\sum_i v_i x_i}_{\text{objective}}
        + \underbrace{\lambda \sum_k P_k(x)^2}_{\text{constraint penalties}}
\end{equation}
where each penalty $P_k(x)^2$ introduces cross-terms between variables.
The number of distinct QUBO variables $q$ determines the Hilbert space
dimension $2^q$ that QAOA must search. We classify problems by how $q$
scales with instance size $n$.

% ── 1.2 Per-problem ──────────────────────────────────────────────────────────
\subsection{Per-Problem Analysis}

\subsubsection*{(a) Knapsack --- Single Capacity Constraint}

\begin{align}
  &x_i \in \{0,1\}, \quad i = 1, \ldots, n \\
  &s \in \{0,1\}^{\lceil \log_2 W \rceil}
    \quad\text{to encode } \textstyle\sum w_i x_i \leq W
\end{align}
\begin{equation}
  \boxed{q_{\text{Knapsack}} = n + \lceil \log_2 W \rceil = O(n + \log W)}
\end{equation}
For $W = O(\mathrm{poly}(n))$: $q = O(n)$.

\textbf{LP boundary structure.} By the simplex vertex theorem applied to a
single-hyperplane feasible region, the LP relaxation has \emph{at most one}
fractional variable at any optimal vertex~\cite{dantzig1957,kellerer2004}.
Hence the boundary zone size $m = O(1)$ and the quantum sub-problem operates
on $O(1)$ qubits.

\subsubsection*{(b) Multi-Dimensional Knapsack --- $k$ Capacity Constraints}

\begin{equation}
  \sum_{j=1}^{n} w_{ij} x_j \leq W_i, \quad i = 1, \ldots, k
\end{equation}
Each inequality requires $\lceil \log_2 W_i \rceil$ slack qubits:
\begin{equation}
  q_{\text{MKP}} = n + k \cdot \lceil \log_2 W_{\max} \rceil = O(n + k \log W)
\end{equation}
For fixed $k$: $O(n)$. For $k = O(n)$: $O(n \log n)$.
LP boundary: up to $\min(k, n)$ fractional variables~\cite{schrijver1986}.

\subsubsection*{(c) Portfolio Optimisation --- Sparse Covariance}

\begin{equation}
  \text{Minimise } -\mu^T x + \lambda\, x^T \Sigma x, \quad x_i \in \{0,1\}
\end{equation}
The objective is already quadratic. For a covariance matrix $\Sigma$ with
$s$ non-zeros per row:
\begin{equation}
  q_{\text{Portfolio}} = n, \quad \text{QUBO terms: } O(ns)
\end{equation}
\begin{equation}
  \boxed{q_{\text{Portfolio}} = O(n)}
\end{equation}
For sparse $\Sigma$ ($s = O(1)$): $O(n)$ terms.
For dense $\Sigma$ ($s = n$): $O(n^2)$ terms but still $n$ variables.

\subsubsection*{(d) Travelling Salesman Problem (TSP)}

Standard position-indexed formulation: $x_{ij} \in \{0,1\}$ (city $i$ at
position $j$), with $n^2$ arc variables plus MTZ position encoding:
\begin{equation}
  q_{\text{TSP}} = n^2 + n \lceil \log_2 n \rceil = O(n^2)
\end{equation}
The fraction of valid TSP tours in the full Hilbert space vanishes doubly
exponentially:
\begin{equation}
  \frac{(n-1)!/2}{2^{n^2}}
  \sim \frac{\sqrt{2\pi n}(n/e)^n}{2 \cdot 2^{n^2}} \to 0
\end{equation}
\begin{equation}
  \boxed{q_{\text{TSP}} = O(n^2), \quad
         \text{valid fraction} = O\!\left(\frac{n!}{2^{n^2}}\right)}
\end{equation}

\subsubsection*{(e) VRP --- $n$ Nodes, $m$ Vehicles (MTZ Formulation)}

Three-index arc variables $x_{ijk} \in \{0,1\}$ (vehicle $k$ on arc
$(i,j)$); MTZ position variables $u_{ik}$; capacity slack:
\begin{equation}
  q_{\text{VRP}} = n^2 m + nm\lceil\log_2 n\rceil
                 + m\lceil\log_2(Q/d_{\min})\rceil = O(n^2 m)
\end{equation}
\emph{Phase 2 validation}: $n=20$, $m=4$ gives
$1680 + 400 + 12 = 2092$ logical qubits (Phase 2f reported 2,648;
27\% difference from T-gate ancilla).
\begin{equation}
  \boxed{q_{\text{VRP}} = O(n^2 m)}
\end{equation}

\subsubsection*{(f) MaxCut --- $m$ Edges, $n$ Nodes}

\begin{equation}
  H_{\text{MaxCut}} = \sum_{(i,j) \in E} \left(x_i + x_j - 2x_i x_j\right),
  \quad x_i \in \{0,1\}
\end{equation}
No feasibility constraints: every $x \in \{0,1\}^n$ is valid.
\begin{equation}
  \boxed{q_{\text{MaxCut}} = n, \quad \text{valid fraction} = 1.0}
\end{equation}

\subsubsection*{(g) Graph Colouring --- $n$ Nodes, $k$ Colours}

$x_{ic} \in \{0,1\}$ (node $i$ receives colour $c$); $n \times k$ variables:
\begin{equation}
  q_{\text{Colouring}} = nk + \text{slack} \approx nk
\end{equation}
\begin{equation}
  \rho = \frac{|E|\, k}{\binom{nk}{2}} \approx \frac{2|E|}{n^2 k}
\end{equation}
\begin{equation}
  \boxed{q_{\text{Colouring}} = O(nk)}
\end{equation}

\subsection{Summary Table}

\begin{table}[ht]
\centering
\caption{QUBO encoding complexity across seven problem classes.}
\label{tab:summary}
\begin{tabular}{@{}lllll@{}}
\toprule
Problem & QUBO vars $q$ & Growth & Constraint density & Valid fraction \\
\midrule
Knapsack         & $n + O(\log W)$ & $O(n)$     & $O(1/n)$      & $O(1)$ \\
Multi-KP (fixed $k$) & $n + k\log W$ & $O(n)$ & $O(k/n^2)$    & $O(1/k)$ \\
Portfolio (sparse) & $n$           & $O(n)$     & $O(s/n)$      & $O(1)$ \\
MaxCut           & $n$             & $O(n)$     & 0 (none)      & $\mathbf{1.0}$ \\
Graph Colouring  & $nk$            & $O(nk)$    & $O(1/k)$      & $O((k!/k^n)^n)$ \\
TSP              & $n^2$           & $O(n^2)$   & $O(1/n^2)$    & $O(n!/2^{n^2})$ \\
VRP              & $n^2 m$         & $O(n^2 m)$ & $O(1/n^2 m)$  & $O(n!^m/2^{n^2m})$ \\
\bottomrule
\end{tabular}
\end{table}

\begin{remark}
Problems with $O(n)$ encoding are either unconstrained (MaxCut, Portfolio)
or have a single dominant constraint (Knapsack). Problems with $O(n^2)$
encoding all arise from pairwise coupling of nodes.
\end{remark}

% ══════════════════════════════════════════════════════════════════════════════
\section{The Barren Plateau -- Encoding Complexity Link}
\label{sec:barren}

\subsection{The Gradient Variance Theorem}

The following is proven in~\cite{mcclean2018barren}:

\begin{theorem}[McClean et al.\ 2018]
For a variational quantum circuit $U(\theta)$ that forms an approximate
unitary 2-design on $n$ qubits, the variance of the gradient of the expected
energy satisfies:
\begin{equation}
  \Var\!\left[\frac{\partial \langle H \rangle}{\partial \theta_k}\right]
  = O\!\left(\frac{1}{2^n}\right)
\end{equation}
\end{theorem}

\noindent This means: to measure the gradient with precision $\varepsilon$,
one needs $O(2^n / \varepsilon^2)$ circuit evaluations.

\subsection{Encoding Size Amplifies the Plateau}

\begin{conjecture}[C1 -- Encoding-Amplified Barren Plateau]
\label{conj:C1}
If a problem requires $q(n)$ QUBO variables, the effective QAOA circuit
operates on $q(n)$ qubits. Therefore:
\begin{equation}
  \Var\!\left[\frac{\partial \langle H \rangle}{\partial \theta_k}\right]
  = O\!\left(\frac{1}{2^{q(n)}}\right)
\end{equation}
\end{conjecture}

\noindent (Formal proof would require extending McClean et al.\ to structured
QUBO circuits.) For VRP with $n=20$, $m=4$: the gradient variance is
$O(2^{-2092}) \approx 10^{-630}$ --- numerically indistinguishable from zero.

\begin{table}[ht]
\centering
\caption{Barren plateau severity by encoding class.}
\label{tab:barren}
\begin{tabular}{@{}llll@{}}
\toprule
Encoding & Effective qubits & Gradient variance
  & Evaluations for $\varepsilon$-gradient \\
\midrule
$O(n)$     & $n$     & $O(2^{-n})$     & $O(2^n/\varepsilon^2)$ \\
$O(n^2)$   & $n^2$   & $O(2^{-n^2})$   & $O(2^{n^2}/\varepsilon^2)$ \\
$O(n^2 m)$ & $n^2 m$ & $O(2^{-n^2 m})$ & doubly-exponential \\
\bottomrule
\end{tabular}
\end{table}

\subsection{Can $O(n)$ Encoding Avoid Barren Plateau?}

\begin{conjecture}[C2 -- Three-Condition Quantum Advantage Criterion]
\label{conj:C2}
Quantum advantage via QAOA is theoretically accessible only when all three
of the following hold simultaneously:
\begin{equation}
  \begin{cases}
    q(n) = O(n) & \text{(linear encoding)} \\
    \text{cost function local} & \text{($O(1)$ qubit support per term)} \\
    \text{valid fraction} \approx 1 & \text{(no feasibility penalty)}
  \end{cases}
\end{equation}
\end{conjecture}

\noindent Evidence: Cerezo et al.~\cite{cerezo2021cost} prove that local cost
functions avoid exponential gradient decay in shallow circuits, while global
cost functions suffer barren plateaus even at depth $O(1)$.
For MaxCut: each edge term $(x_i + x_j - 2x_i x_j)$ acts on 2 qubits
($\Rightarrow$ local). For Knapsack: the penalty $\lambda(\sum w_i x_i - W)^2$
involves all $n$ variables simultaneously ($\Rightarrow$ global, full barren
plateau despite $O(n)$ encoding).

\noindent MaxCut satisfies all three conditions. Knapsack satisfies only the
first. VRP satisfies none.

% ══════════════════════════════════════════════════════════════════════════════
\section{The Circularity Conjecture}
\label{sec:circularity}

\subsection{Formal Statement}

\begin{definition}[Problem Sets $\calA$, $\calB$, $\calC$]
\label{def:ABC}
\begin{align}
  \calA &= \{\text{problems where LP boundary zone} = O(1)\} \\
  \calB &= \{\text{problems where QUBO encoding} = O(n)\} \\
  \calC &= \{\text{problems where LP + rounding achieves }
              (1-\varepsilon)\text{-approximation}\}
\end{align}
\end{definition}

\begin{conjecture}[C3 -- Circularity Theorem]
\label{conj:C3}
For natural combinatorial optimisation problem classes:
\begin{equation}
  \calA \subseteq \calB, \quad \calA \subseteq \calC,
  \quad\text{and}\quad \calB \cap \calC \subseteq \calA
\end{equation}
\emph{Informally:} The problems where quantum sub-problems are small enough
to be tractable are exactly the problems where classical LP already works well.
\end{conjecture}

\subsection{The $\calA \subseteq \calB$ and $\calA \subseteq \calC$ Directions}

These follow from LP integrality theory~\cite{schrijver1986}.
A system $Ax \leq b$ is \emph{Totally Dual Integral (TDI)} if for every
integer objective $c$, the LP dual has an integer optimal solution.
If additionally $b$ is integer, the LP has an integer optimal solution
(integrality gap = 1). A sufficient condition: $A$ is \emph{totally
unimodular (TU)} (every square sub-determinant $\in \{-1,0,+1\}$).

\noindent\textbf{The $\calA \subseteq \calB$ argument:}
\begin{align}
  &\text{Small LP boundary}
    \;\xLeftrightarrow{\text{structural}}\;
    \text{LP relaxation near-integral}
    \;\xLeftrightarrow{\text{theory}}\;
    \text{small integrality gap} \notag\\
  &\text{Small LP boundary}
    \;\Rightarrow\;
    \text{QUBO encodes only boundary}
    \;\Rightarrow\;
    O(1)\text{ extra variables}
    \;\Rightarrow\;
    \calA \subseteq \calB \quad\textit{[Conjecture]}
\end{align}

\noindent\textbf{The $\calA \subseteq \calC$ argument:}
\begin{equation}
  \text{Small integrality gap}
  \;\Rightarrow\;
  \text{LP + rounding achieves }(1-\varepsilon)\text{-approximation}
  \;\Rightarrow\;
  \calA \subseteq \calC
\end{equation}
For Knapsack this is the FPTAS argument~\cite{vazirani2001}.

\subsection{The Reverse Direction: $\calB \cap \calC \subseteq \calA$}

\begin{openprob}[B$\cap$C$\subseteq$A Direction]
\label{op:BcapC}
If a problem has $O(n)$ QUBO encoding \emph{and} LP + rounding works well,
does its LP relaxation boundary zone necessarily satisfy $m = O(1)$?
\end{openprob}

\noindent\textbf{Informal argument (not a proof).}
Suppose $P \in \calB \cap \calC$. Since $P \in \calB$, the QUBO has $O(n)$
variables, which requires each penalty $P_k$ to have $O(1)$ variable support.
Since $P \in \calC$, LP + rounding succeeds, implying a bounded integrality
gap. By the Hoffman--Kruskal theorem~\cite{hoffmankruskal1956,schrijver1986},
a bounded gap implies the LP optimal face has bounded dimension, hence bounded
fractional variables, hence $P \in \calA$.

\noindent\textbf{Where the argument fails.}
The step ``bounded integrality gap $\Rightarrow$ bounded LP face dimension''
is not a proven theorem in general. It holds for TU constraint matrices but
may fail for general $\{0,1\}$-programs with bounded gap.

\noindent\textbf{Three routes to a formal proof.}
\begin{enumerate}
  \item \emph{TDI + Hoffman--Kruskal:} If $Ax \leq b$ is TDI with sparse
    $A$ (each row $O(1)$ non-zeros), then the LP optimal face has integral
    vertices. Gap: sparsity of $A$ must be connected to $O(n)$ QUBO
    encoding~\cite{hoffmankruskal1956,schrijver1986}.
  \item \emph{Communication complexity:} Use Yao's minimax
    principle~\cite{yao1979,kushilevitz1997} to show $O(n)$ QUBO encoding
    implies $O(n)$-bit constraint certificates, limiting boundary zone.
    Gap: formalising ``QUBO encodes the problem'' precisely.
  \item \emph{Constraint graph treewidth:} If the constraint graph $G_P$
    is sparse ($O(n)$ edges), the LP feasible region has bounded treewidth,
    limiting fractionality~\cite{bodlaender1993,cygan2015}. Gap: treewidth
    $k$ implies bounded fractionality only for $k = O(1)$.
\end{enumerate}

\noindent\textbf{Verdict.} $\calB \cap \calC \subseteq \calA$ is a
well-motivated open problem with three plausible proof routes. Proving it
would constitute a significant result connecting LP integrality theory, QUBO
encoding, and quantum advantage.

\subsection{Counter-Direction: Large Gap $\Rightarrow$ Large QUBO}

For routing problems (TSP, VRP), achieving even the LP lower bound requires
exponentially many Dantzig--Fulkerson--Johnson subtour elimination
constraints. Encoding these as QUBO penalties forces $O(n^2)$ arc variables.
The QUBO must encode what the LP constraint system requires, providing
structural reason to believe $\calA \cap \calC^c = \emptyset$ for natural
classes.

% ══════════════════════════════════════════════════════════════════════════════
\section{Candidate Exceptions}
\label{sec:exceptions}

\subsection{MaxCut as Primary Candidate}

MaxCut satisfies all three conditions in Conjecture~\ref{conj:C2}: $O(n)$
encoding, all solutions valid, local cost function. The key question is
whether QAOA can beat the Goemans--Williamson 0.878 approximation guarantee.

If the Unique Games Conjecture (UGC) holds, achieving ratio
$> 0.878 + \varepsilon$ for MaxCut is NP-hard
classically~\cite{khot2007maxcut}.

\begin{conjecture}[C4 -- MaxCut Quantum Advantage]
\label{conj:C4}
If UGC is true, then quantum algorithms with polynomial overhead could achieve
$> 0.878$ approximation for MaxCut, representing genuine quantum advantage.
\end{conjecture}

\noindent\textbf{Current empirical status.} QAOA $p=1$ achieves 0.6924
$<$ 0.878. Bravyi et al.~\cite{bravyi2020obstacles} showed QAOA $p < n/2$
cannot beat Goemans--Williamson on certain graph families.
Farhi et al.~\cite{farhi2020typical,farhi2020worst} showed that for depth
$p < c \log n$, the QAOA is limited by locality barriers on random
$d$-regular graphs. We include Conjecture~\ref{conj:C4} as a candidate exception to the Circularity Conjecture; we do not prove it, and its resolution depends on the status of the Unique Games Conjecture.

\subsection{Quantum Local Hamiltonian ($k$-Local Hamiltonian)}

The $k$-Local Hamiltonian problem is QMA-complete~\cite{kitaev2002} for
$k \geq 2$ --- the quantum analogue of NP-completeness. Its natural encoding
is $O(n)$ qubits with $O(n)$ local terms. Problems with genuine quantum
structure (where cost functions involve quantum superposition) have no
classical approximation barrier analogous to the UGC barrier for classical
combinatorial optimisation. VQE for quantum chemistry is in this class.

\subsection{QUBO at Phase Transitions}

\begin{proposition}[\emph{(empirical)} -- Phase Transition Hardness]
\label{prop:phase}
Random $k$-SAT near the satisfiability phase transition ($\alpha = \alpha_c$)
is computationally hard for classical solvers, with exponentially many
near-degenerate local minima. Whether quantum tunnelling provides an advantage
is an open experimental question. Hastings~\cite{hastings2021} showed that
for random MAX-2-SAT, simulated annealing and QAOA have similar performance
at shallow depths, suggesting phase transition hardness does not automatically
grant quantum advantage.
\end{proposition}

% ══════════════════════════════════════════════════════════════════════════════
\section{Implications for the Phase 2 VRP Results}
\label{sec:vrp}

\subsection{Consistency Check: 1.2M Physical Qubits}

The Phase 2 benchmark (20-ATM CVRP, 4 vehicles) reported
$\approx 1{,}188{,}952$ physical qubits. Our theoretical prediction:
\begin{align}
  q_{\text{VRP}}(n=20, m=4)
    &= n^2 m + nm\lceil\log_2 n\rceil + m\lceil\log_2(Q/d_{\min})\rceil \\
    &= 400 \times 4 + 20 \times 4 \times 5 + 4 \times 3
     = 1680 + 400 + 12 = 2092\text{ logical}
\end{align}
With surface code distance $d=15$: $1\text{ logical qubit} \rightarrow 2d^2-1=449$
physical:
\begin{equation}
  2092 \times 449 = 939{,}308\text{ physical}
\end{equation}
The 27\% gap (939K vs.\ 1.19M) arises from T-gate ancilla in magic state
distillation for fault-tolerant $R_z$ gates. This is consistent with the
theoretical $O(n^2m)$ class.

\subsection{Alternative Formulations and Encoding Lower Bounds}

The Dantzig--Fulkerson--Johnson (DFJ) flow formulation reduces to
$q_{\text{DFJ}} = n^2\lceil\log_2(m+1)\rceil = O(n^2 \log m)$;
for $n=20$, $m=4$: 1200 logical, 538,800 physical (55\% reduction), but
still in the $O(n^2)$ class. The quadratic bottleneck remains.

\textbf{Assessment of the subtour encoding lower bound.}
A counting argument over $2^n - 2$ proper subsets yields only
$\lceil\log_2(2^n)\rceil = n$ bits --- an $\Omega(n)$ lower bound,
\emph{not} $\Omega(n\log n)$.

\begin{openprob}[Subtour Encoding Lower Bound -- Conjecture C5]
\label{op:C5}
The minimum QUBO variable count for any correct encoding of TSP/VRP subtour
constraints satisfies:
\begin{equation}
  q_{\min} = \Omega(n \log n)
\end{equation}
\end{openprob}

\noindent\textbf{Informal supporting argument (proof sketch, not theorem).}
A TSP optimal tour is a permutation $\pi \in S_n$; there are $n!$ distinct
tours. By Stirling's approximation:
\begin{equation}
  \lceil\log_2(n!)\rceil = \Theta(n \log n)
\end{equation}
Any encoding distinguishing all optimal tours needs $\Omega(n\log n)$ bits.
\emph{Gap:} this bounds the solution encoding, not the constraint encoding.
Multiple QUBO assignments can map to the same tour, and the two need not
coincide. A formal proof requires connecting information-theoretic solution
lower bounds to constraint lower bounds --- currently open.

\subsection{The Bug-Lessons as Theoretical Signals}

Three bugs discovered in the Phase 2 QAOA implementation encode theoretical
content:

\begin{table}[ht]
\centering
\caption{Phase 2 implementation bugs as theoretical signals.}
\label{tab:bugs}
\begin{tabular}{@{}lll@{}}
\toprule
Bug & Root Cause & Theoretical Implication \\
\midrule
$\lambda$ too small &
  Penalty $\ll$ objective &
  $\lambda \geq \mathrm{Lip}(f)/\text{constraint curvature}$ required \\
MILP infeasible &
  Over-constrained model &
  Constraint satisfaction in NP; tight encoding risks infeasibility \\
$\langle H\rangle$ vs.\ argmax &
  Two distinct quantities &
  Expectation is quantum observable; argmax is classical post-processing \\
\bottomrule
\end{tabular}
\end{table}

% ══════════════════════════════════════════════════════════════════════════════
\section{Open Questions}
\label{sec:open}

\subsection*{Q1: Can the Circularity Conjecture Be Proven?}

A complete proof requires three steps:
\begin{enumerate}
\item LP boundary zone $O(1)$ $\Rightarrow$ QUBO encoding $O(n)$
  (from LP vertex degree-of-freedom counting).
\item $O(n)$ QUBO + LP integrality gap $< \varepsilon$
  $\Rightarrow$ $(1-\delta(\varepsilon))$-approximation
  (FPTAS argument~\cite{vazirani2001}).
\item LP + rounding fails $\Rightarrow$ gap $> 1$
  $\Rightarrow$ $> O(n)$ QUBO
  (requires connecting LP theory with combinatorial topology).
\end{enumerate}

\subsection*{Q2: Which Problem Class Is Closest to Quantum Advantage?}

\begin{enumerate}
\item \textbf{$k$-Local Hamiltonian (QMA-complete):} Quantum advantage
  proven complexity-theoretically --- but not a classical combinatorial problem.
\item \textbf{MaxCut on expander graphs:} Theoretically possible if QAOA
  depth $p = \Omega(n)$ and UGC holds; requires fault-tolerant hardware.
\item \textbf{Quadratic Assignment Problem (QAP) at phase transition:}
  No known efficient classical algorithm; QUBO encoding $O(n^2)$ --- barren
  plateau challenge remains.
\item \textbf{Protein folding (HP model):} $O(n)$ encoding on a lattice,
  quantum-native structure.
\item \textbf{Knapsack/VRP/TSP:} Structural barriers make quantum advantage
  unlikely without encoding breakthroughs.
\end{enumerate}

\subsection*{Q3: Minimum Structure for Meaningful Hybrid?}

Based on this analysis, the minimum conditions for a meaningful hybrid
quantum-classical approach are:
\begin{equation}
  \begin{cases}
    1. & \text{QUBO encoding: } O(n) \text{ variables} \\
    2. & \text{Cost function: local ($O(1)$ qubit support per term)} \\
    3. & \text{Valid solution fraction: bounded away from 0} \\
    4. & \text{Classical approximation: below UGC barrier} \\
    5. & \text{LP boundary zone: } \Omega(\log n) \text{ (non-trivial)}
  \end{cases}
\end{equation}
No known classical combinatorial optimisation problem satisfies all five
simultaneously at practical scale. The most promising research direction:
\emph{problem-specific mixers} (XQAOA, ADAPT-VQE) that preserve feasibility
constraints in superposition.

% ══════════════════════════════════════════════════════════════════════════════
\section{Related Work}
\label{sec:related}

\textbf{Bravyi, Kliesch, Koenig, Tang (2022)~\cite{bravyi2022hybrid}.}
This paper analyses Recursive QAOA (RQAOA) on MAX-$k$-CUT, equivalent to
approximate $k$-colouring. Standard (non-recursive) QAOA fails for most
regular bipartite graphs at any constant depth $p$, and level-1 RQAOA
admits an efficient classical simulation. Relation to this work: Graph
Colouring has $O(nk)$ QUBO encoding (Table~\ref{tab:summary}), placing it
in the super-linear class. The result is consistent with
Conjecture~\ref{conj:C1}: barren plateau pressure scales as $O(2^{-nk})$.
Additionally, the ``one colour per node'' constraint is global, violating
the locality condition of Conjecture~\ref{conj:C2}.

\textbf{Farhi, Gamarnik, Gutmann (2020)~\cite{farhi2020typical,farhi2020worst}.}
For depth $p < c\log n$, QAOA on Maximum Independent Set on random
$d$-regular graphs cannot achieve better than approximately 0.854 times
the optimum. The mechanism is locality: shallow circuits see only tree-like
$p$-hop neighbourhoods. This is the strongest evidence for the locality
condition in Conjecture~\ref{conj:C2}. The depth requirement $p = \Omega(\log n)$
places QAOA precisely in the regime where barren plateaus are catastrophic
--- squeezing quantum advantage from both sides.

\textbf{Cerezo et al.\ (2021)~\cite{cerezo2021cost}.}
Cost function dependent barren plateaus: local cost functions (observable on
$O(1)$ qubits) have at most polynomially decaying gradients for shallow
circuits, while global cost functions suffer exponential decay even at depth
$O(1)$. This closes the gap in Conjecture~\ref{conj:C2}: $O(n)$ encoding
is necessary but not sufficient --- locality of the cost function is an
independent requirement.

% ══════════════════════════════════════════════════════════════════════════════
\section*{Conclusion}

Classical combinatorial optimisation problems are classical by definition:
their solutions are binary strings. Quantum advantage for such problems
requires outperforming the best classical approximation. For problems with
$O(n^2)$ or higher QUBO encoding, barren plateau gradient suppression scales
as $O(2^{-n^2})$ --- numerically zero for any practical $n$. The problems
where quantum circuits are small enough to avoid this barrier are exactly
those where classical LP already finds near-optimal solutions. This
circularity is structural, not a hardware limitation, and will not be
resolved by improved quantum chips alone.

This work does not claim quantum computing is without promise. It identifies
that the right problem class for gate-based QAOA applied to classical
combinatorial optimisation has not yet been found, and provides a structural
framework for the search.

% ── Data availability ─────────────────────────────────────────────────────────
\section*{Data and Code Availability}

All code, data, and analysis scripts are available at\\
\url{https://github.com/jasstt/quantum-knapsack-vrp-theory}\\
under the MIT License. The predecessor project (20-ATM CVRP empirical data)
is at \url{https://github.com/jasstt/stochastic-VRP-decision-focused-learning}.

% ── References ────────────────────────────────────────────────────────────────
\bibliographystyle{plain}
\bibliography{references}

\end{document}
